\section{Introdução}

A inteligência artificial (IA) e o aprendizado de máquina (\textit{machine learning} --- ML) constituem o núcleo cognitivo e o diferencial axiomático dos gêmeos digitais odontológicos contemporâneos. Enquanto os capítulos anteriores estabeleceram as rigorosas bases da modelagem geométrica e biomecânica (\autoref{ch:fundamentos}) e o panorama evolutivo da odontologia digital (\autoref{ch:panorama}), o presente capítulo aprofunda os algoritmos, as arquiteturas profundas e os frameworks de orquestração que conferem aos gêmeos digitais sua capacidade de aprendizado contínuo, predição antecipatória e otimização autônoma. 

Do ponto de vista epistemológico, um gêmeo digital odontológico desprovido de IA reduzir-se-ia a um modelo tridimensional inerte --- uma representação visual topologicamente precisa, porém desprovida de agência e inteligência analítica. A transição paradigmática de representações estáticas para ontologias dinâmicas ocorre exclusivamente pela integração de camadas de IA. É esse arcabouço algorítmico que permite ao gêmeo digital segmentar automaticamente estruturas anatômicas complexas em exames de imagem volumétricos, prever com precisão submilimétrica a trajetória de movimentação dentária sob estímulos ortodônticos, otimizar o posicionamento de fixações implantares com base em simulações biomecânicas multifísicas e detectar anomalias incipientes que sistematicamente escapariam ao escrutínio humano \cite{duggal2026, katsoulakis2024}.

A literatura científica aponta uma aceleração notável no campo da IA odontológica a partir do biênio 2024--2026, impulsionada por três vetores de força convergentes: (1) a consolidação de \textit{datasets} clínicos digitalizados e multimodais; (2) o avanço estrutural de arquiteturas de aprendizado profundo (\textit{deep learning}); e (3) a integração de plataformas de orquestração multiagente, capazes de modular tarefas complexas \cite{khoshfekr2025, menon2026}. Como observam Roy et al. (2025), a integração destas inovações aos gêmeos digitais odontológicos transcende uma mera evolução incremental, consubstanciando uma verdadeira quebra de paradigma na semiologia e na terapêutica \cite{roy2025editorial}.

\destaque{
A ascensão de ecossistemas abertos, como o OpenCode Ecosystem, evidencia o movimento irreversível em direção à interoperabilidade e padronização. Ao orquestrar \textit{pipelines} de simulação com validação rigorosa (TDD/SDD), a integração clínica atinge patamares de reprodutibilidade e confiança (TrustEngine) exigidos para a saúde humana.
}

Este capítulo encontra-se estruturado em oito seções complementares. As seções \ref{sec:ia-fundamentos} e \ref{sec:ia-aplicacoes} assentam os pressupostos teóricos (matemáticos e arquiteturais) e as consequentes aplicações práticas da IA em gêmeos digitais odontológicos. As seções subsequentes versam sobre eixos críticos da literatura moderna: IA explicável e governança algorítmica (\ref{sec:ia-xai}), convergência com modelos de linguagem de grande escala (\ref{sec:ia-llms}), descoberta de biomarcadores via aprendizado profundo (\ref{sec:ia-biomarcadores}) e, de forma imperativa, as barreiras de implementação (\ref{sec:ia-desafios}) e horizontes prospectivos (\ref{sec:ia-conclusão}) de um setor projetado para transformar diametralmente a saúde pública e privada.
